\begin{abstract}
Los call centers procesan miles de interacciones diarias en las que episodios de ira,
frustración o ansiedad pueden derivar en escalamientos que deterioran la calidad del
servicio y el bienestar de los agentes. La evaluación manual de calidad solo cubre el
1--5\,\% de las llamadas, dejando sin monitorear la gran mayoría de los momentos
críticos. Este trabajo propone un sistema automatizado de detección de picos
emocionales basado en el análisis de características paralingüísticas —como
\textit{pitch} (frecuencia fundamental F0), intensidad, jitter, shimmer y MFCCs
(coeficientes cepstrales en las frecuencias de Mel)—, combinado con algoritmos de
machine learning (Z-scores estadísticos, Isolation Forest, One-Class SVM y Random
Forest). Se evalúan tres métodos para el cálculo de la línea base emocional
individualizada por hablante. El sistema apunta a una precisión $\geq 75$\,\% y un
recall $\geq 70$\,\% en llamadas reales en español, con validación mediante anotaciones
de expertos.
\end{abstract}

\begin{IEEEkeywords}
reconocimiento de emociones en voz, análisis paralingüístico, call center, detección
de anomalías, machine learning, línea base emocional, Isolation Forest
\end{IEEEkeywords}
